\subsection{The original coefficient actions on coordinate sections}

\begin{lemma}\label{lem:original-native-coordinate-coefficients}
\lean{EnvelopeProperty.ProjBerk.nativeCoordinateSectionsIso_coefficients}
\leanok
\uses{projective:native-scheme-coordinate-sections,projective:native-structure-morphism,projective:native-relative-coordinate-chart,scheme:relative-section-coefficients}
Let $R$ be any commutative ring, let $P=\mathbf P^n_R$ be the original
native polynomial projective scheme, let $g:P\to\operatorname{Spec}R$
be its original structure map, and let $U_j=D_+(x_j)$.
For the existing whole coordinate-section isomorphism
$e_j:B_j=R[u_i:i\ne j]\simeq\Gamma(P,U_j)$ and every $r\in R$,
\[
 e_j(C(r))=\operatorname{sectionBaseMap}(g,U_j)(r).
\]
The right side is the already defined coefficient map obtained by pulling
back the original base global section and restricting to $U_j$.
\end{lemma}
\begin{proof}
Let $h:U_j\simeq\operatorname{Spec}B_j$ be the existing actual coordinate
scheme isomorphism and $u:U_j\to P$ its original open inclusion. The
existing coefficient triangle is
$h^{-1}u g=\operatorname{Spec}(C:R\to B_j)$.
The existing description of $e_j$ is the composite of the inverse
global-section isomorphism of $\operatorname{Spec}B_j$, pullback along
$h$, and the actual identification of the open's global sections with
$\Gamma(P,U_j)$. Apply global sections contravariantly to this coefficient
triangle, and use the naturality of the global-section isomorphisms of
affine spectra. It says that applying $e_j^{-1}$ to the original coefficient
pullback gives exactly $C(r)$. Restricting the pullback along $g$ to $U_j$
is the pullback along $u g$, by the original sheaf restriction naturality;
this is exactly the definition of $\operatorname{sectionBaseMap}(g,U_j)$.
Thus $e_j^{-1}(\operatorname{sectionBaseMap}(g,U_j)(r))=C(r)$.
Applying the inverse equivalence $e_j$ proves the assertion. No equality
of sections is deduced merely from equality at points.
\end{proof}

\begin{lemma}\label{lem:original-native-quotient-coordinate-algebra}
\lean{EnvelopeProperty.ProjBerk.nativeQuotientCoordinateSectionsAlgEquiv}
\leanok
\uses{lem:original-native-coordinate-coefficients,lem:original-native-quotient-coordinate-sections,scheme:actual-section-coefficient-naturality,scheme:relative-section-coefficients}
Let $J$ be any original homogeneous ideal, let $i:Z\to P$ be the actual
inclusion of its original kernel closed subscheme, let $V_j=i^{-1}(U_j)$,
and let $D_j\subset B_j$ be its original dehomogenized ideal.
Give $\Gamma(Z,V_j)$ the already defined coefficient action
$\operatorname{sectionBaseAlgebra}(i g,V_j)$, coming from its original
structure morphism, and give $B_j/D_j$ its original quotient $R$-algebra
structure. The previously constructed whole-ring isomorphism $\Phi_j$
is an $R$-algebra isomorphism for precisely these actions.
Its underlying ring equivalence is exactly the original $\Phi_j$.
No new coefficient action, primeness, field, reducedness or relevance
assumption is introduced.
\end{lemma}
\begin{proof}
Use the existing whole section-coefficient naturality theorem with the
actual inclusion $i$, the two actual structure morphisms, and the identity
homomorphism of $R$. The coefficient triangle is the defining equality
$i g=i g$. For every $r$, it gives
\[
 i^*(\operatorname{sectionBaseMap}(g,U_j)(r))
   =\operatorname{sectionBaseMap}(i g,V_j)(r).
\]
The preceding lemma identifies the left input with $e_j(C(r))$.
The existing every-original-polynomial formula for $\Phi_j$, applied to
$p=C(r)$, now gives
\[
 \Phi_j(\operatorname{sectionBaseMap}(i g,V_j)(r))=[C(r)].
\]
The source coefficient map is exactly the algebra map of the existing
named section action; the target coefficient map is exactly $r\mapsto[C(r)]$.
Thus the already constructed original ring equivalence commutes with both
original coefficient maps. Equip that same equivalence with this proved
commutation law to obtain the claimed algebra isomorphism. Its inverse
is the original ring inverse, so no carrier, operations or sections change.
\end{proof}
